\section{MCP Servers：工具扩展}

\subsection{MCP 的角色}

MCP（Model Context Protocol）servers 主要用于将工具（tools）接入 coding agent，扩展其文件 I/O 和 bash 命令之外的能力。MCP 规范还包括 resources、prompts、elicitations 等特性，但这些在 MCP 客户端和 coding agent harness 中通常支持不佳。

当你将 MCP server 接入 coding agent 时，可用工具列表、工具描述和调用参数会被注入 agent 的 system prompt。因此，MCP server 可以通过工具描述来定制 agent 的行为。

\begin{warningbox}{MCP 安全警告}
因为 MCP server 的工具描述会被添加到 coding agent 的 system prompt 中，\textbf{永远不要连接你不信任的 MCP server}。这可能成为 prompt injection 的危险载体。STDIO servers 和其他通过 \texttt{npx} 或 \texttt{uvx} 在客户端运行的 server 也可以在没有 prompt injection 的情况下在你的主机上执行代码。
\end{warningbox}

\subsection{工具过多的危害}

HumanLayer 的亲身经验：接入过多 MCP 工具后，context window 被工具描述填满，agent 更快地进入``笨区''（dumb zone）。

\begin{knowledgebox}{Instruction Budget 概念}
每个 agent 的 context window 有一个隐式的 ``instruction budget''（指令预算）。每一条无关的工具描述都是 agent 必须处理但没有收益的指令。这些无用指令会挤占真正重要的上下文空间。Anthropic 甚至为此发布了实验性的 MCP tool search 功能，在用户连接了过多 MCP 工具时渐进式地向 Claude 披露工具。
\end{knowledgebox}

\textbf{关键建议}：如果你不在主动使用某个提供大量工具的 MCP server，就关掉它。

\subsection{CLI 优于 MCP 的场景}

HumanLayer 发现，当 MCP server 复制了训练数据中已有 CLI 的功能时，直接提示 agent 使用 CLI 效果更好。对于 GitHub、Docker 或大多数数据库，coding agent 可以直接使用对应的 CLI 和 shell 命令。模型在训练时已经见过这些工具，且 CLI 可以与 \texttt{grep}、\texttt{jq} 等工具组合使用，提供额外的上下文效率。

\subsection{实战案例：Linear CLI 替代 MCP}

HumanLayer 曾使用 Linear MCP server，后来发现实际只用到其中一小部分功能。于是他们编写了一个轻量 CLI 来封装 Linear API，提供非常精简的响应，并在 CLAUDE.md 中给出 6 条用法示例。

\begin{lstlisting}[caption={CLAUDE.md 中的 Linear CLI 示例配置},language={}]
## Linear
Use the Linear CLI for:
- fetching issues:    linear get-issue ENG-XXXX
- listing issues:     linear list-issues or linear my-issues
- adding comments:    linear add-comment -i ENG-XXXX "comment"
- adding links:       linear add-link ENG-XXXX "url" -t "title"
- updating status:    linear update-status ENG-XXXX "status name"
- get branch name:    linear get-issue-v2 ENG-XXXX --fields branch
- get ticket images:  linear fetch-images ENG-XXXX
\end{lstlisting}

这种方式节省了数千个 token（原本来自 MCP server 的工具定义），以及更多来自冗长 MCP 响应的 token。

\subsection{本章小结}

MCP servers 是扩展 agent 能力的强大工具，但需要克制使用。工具过多会挤占 context window、增加无效指令负担。当 CLI 已经在训练数据中被充分覆盖时，CLI 是更好的选择。始终以 context efficiency 为优化目标。


